% Standalone report fragment: requires amsmath, amssymb, amsthm, booktabs.
\section{PA194NS: crossed command semantics and a complete proof compiler}
\label{sec:ns-proof}

PA194NS combines the no-push, top-result modus-ponens command of PA194N
with the cache-only symmetric Shift implementation of PA194CS.
The underlying numerical axiom calculus and its ordered emitter are
unchanged. The combination is not an immediate source equivalence with
either parent: the proof-command compiler must be checked again, because
its old final contradiction detector no longer works.

The machine has 194 working states, 388 rows, start state 0, halt label 194,
a two-way binary tape initially zero, and only left/right moves. In native
labels, the entire edit from PA194CS is
\[
 (164,1):\quad 1R130\ \longmapsto\ 1R114.
\]
The literal state count remains
\[
 135\text{ fixed controls}+67\text{ AX readers}
       -7\text{ AX folds}-1\text{ controller fold}=194.
\]
The seven AX folds and the controller identification $41\mapsto330$ are
verified against the final table simultaneously, with 32 local cases.
The source now has 3,650 primitive sites and 4,750 branch outcomes.

\subsection{Exact command semantics}

Write
\[
 p(a,b)=\frac{(a+b)(a+b+1)}2+b,\qquad [a;w]=p(a,w),
 \qquad I(a,b)=4p(a,b),\qquad U(i,a)=4p(a,i)+2.
\]
A numerical stack has the literal identities
\[
 \operatorname{hd}[a;w]=a,\quad\operatorname{tl}[a;w]=w,
 \quad[0;0]=0,\quad\operatorname{hd}(0)=\operatorname{tl}(0)=0.
\]
Thus the notation does not assume that numerical stacks form a free-list
algebra. Code zero is truth, code one is the designated contradiction,
and $I(0,0)=0$ is an intentional constructor alias.

A command store is $(t,w,d,h)$: current top, Cantor stack, numerical offset,
and cached formula. For $b=d+x$, the new MP and Shift commands are
\begin{align}
 \mathsf{MP}_x(t,w,d,h)&=
 \begin{cases}
 (b,w,b,0),&I(t,b)=h,\\
 (0,w,b,0),&I(t,b)\ne h,
 \end{cases}\label{eq:ns-mp}\\
 S(t,w,d,h)&=
 \begin{cases}
 \mathsf{halt},&h=1,\\
 (h,\operatorname{tl}(w),d,\operatorname{hd}(w)),&h\ne1.
 \end{cases}\label{eq:ns-shift}
\end{align}
The AX command is the unchanged separate-head emitter. For
$D=d+x$ and $H=h+y$,
\begin{equation}
 \mathsf{AX}_{A,x,y}(t,w,d,h)
 =\bigl(g(H,D,A),[q_1,\ldots,q_9,t;w],0,0\bigr).
 \label{eq:ns-ax}
\end{equation}
Here $q_1=\Gamma(A)=I(A,U(A,A))$, $q_2,\ldots,q_7$ are the
six inherited PA196 heads in reverse emission order, and
\[
 q_8=\rho(D,A,H,d)=I(I(I(H,d),D),I(I(D,H),I(A,H))),
 \qquad q_9=U(d,h).
\]
The top $g$ is the existing guarded separation gate. The executable
free-constructor audit checks every complete formula tree, parameter,
stack position, and guard suffix from the delivered grammar. All
$q_1,\ldots,q_8,g$ have residue zero modulo four; $q_9$ has residue two
and is at least two. None is code one, for any natural input store.
This arithmetic fact requires no consistency assumption.

These semantics distinguish PA194NS from both parents. From
$(t,w,d,h)=(1,0,0,0)$, the PA194N Shift halts, whereas the PA194NS Shift
returns $(0,0,0,0)$. Also, $p(0,1)=2$, so $I(0,1)=8$;
$\mathsf{MP}_1(0,0,0,8)$ is $(1,0,1,0)$ in PA194NS and
$(0,0,0,1)$ in PA194CS. A successful implication therefore lands in a
different register from the register tested for contradiction. The next
plain Shift can discard a false current top; the compiler below uses an
explicit detector that repairs this mismatch.

\subsection{Binding the new source block}

In the PA194CS source, the success block has exactly the three sites
3634--3636, at binary prefix \texttt{1001010111}. It transfers $D$ to $H$.
PA194NS replaces that subtree by the following seven-site block:
\[
\begin{array}{r|l}
3634&\operatorname{dec}D\;[3638,3635]\\
3635&\operatorname{inc}T\to3636\\
3636&\operatorname{inc}X\to3637\\
3637&\operatorname{jump}3634\\
3638&\operatorname{dec}X\;[3641,3639]\\
3639&\operatorname{inc}D\to3640\\
3640&\operatorname{jump}3638.
\end{array}
\]
At successful MP entry, the unchanged preceding constructor/comparison
code supplies $(T,H,X,D)=(0,0,0,b)$. The first loop preserves
$T+D=b$ and $X=T$, and ends with $(T,X,D)=(b,b,0)$.
The second preserves $D+X=b$ and ends with $(T,X,D)=(b,0,b)$.
Both loops terminate by induction on their decremented register, preserve
$W$ and every other register, and perform no stack push. This proves the
success clause of (\ref{eq:ns-mp}); the failure path is unchanged.
The exact source comparison checks every outside instruction after index
relocation and proves that the only outside entrance to the replacement
is the zero exit from site 3628 to site 3634. No outside edge enters its
interior.

The cached symmetric Shift block is retained byte for byte from PA194CS.
It clears $T$, transfers $H$ into $T$ and $X$, increments the now-zero $H$,
and compares $(H,X)=(1,h)$. The inherited destructive comparison halts
exactly on equality and otherwise clears both operands before decoding
one stack cell. Induction on the sum of its two inputs proves these
properties. Thus it realizes (\ref{eq:ns-shift}) under the common command-boundary
invariant $X=0$. With arbitrary incoming scratch $x$, the inherited
comparison of the restored and symmetric tails instead tests $h+x=1$
in both implementations. That stronger tail-equivalence statement is not
being used to assert the command interface for arbitrary scratch.

\subsection{Two AX identities and the changed Pop contract}

Let $A_0=\mathsf{AX}_{0,0,0}$. Directly from
(\ref{eq:ns-ax}) and (\ref{eq:ns-shift}),
\begin{align}
 \mathsf{AX}_{A,x,y};S^9(t,w,d,h)
   &=(q_8,[t;w],0,U(d,h)),\label{eq:ns-nine}\\
 \mathsf{AX}_{A,x,y};S^{10}(t,w,d,h)
   &=(U(d,h),w,0,t).\label{eq:ns-ten}
\end{align}
Every displayed Shift is safe, since its tested cache is zero or one of
the generated non-one codes. In particular these are total identities
for arbitrary natural $t,w,d,h$, including $t=1$.

Define
\[
 \mathsf{Push}=A_0;S^9;\mathsf{MP}_0,\qquad
 \mathsf{Pop}=A_0;S^{11},\qquad
 \mathsf{Discard}=\mathsf{MP}_0.
\]
Push is total and obeys
\[
 \mathsf{Push}(t,w,d,h)=(0,[t;w],0,0).
\]
Indeed its final MP compares an implication of residue zero with the
universal code of residue two in (\ref{eq:ns-nine}), and hence fails with
offset zero. Also, both MP branches coincide when $d=x=0$, giving
\[
 \mathsf{Discard}(t,w,0,h)=(0,w,0,0)
\]
for arbitrary old cache $h$.

The correct Pop contract is now a case distinction:
\begin{equation}
 \mathsf{Pop}(t,[a;w],d,h)=
 \begin{cases}
 \mathsf{halt},&t=1,\\
 (t,w,0,a),&t\ne1.
 \end{cases}\label{eq:ns-pop}
\end{equation}
The first ten Shifts are (\ref{eq:ns-ten}); the eleventh tests exactly the
saved old top $t$. If it passes, Cantor projection gives the stated
store. This applies equally when $a=w=0$, because $[0;0]=0$.
For the zero stack, a nonfalse Pop returns $(t,0,0,0)$; a zero Push onto
the zero stack returns the zero stack. No fictitious sentinel cell is
counted in either identity.

\subsection{Axiom selection and generalization}

For an arbitrary incoming $(t,w,0,0)$ and $1\le j\le7$, the word
\[
 \mathsf{Discard};\mathsf{AX}_{A,D,H};S^{j+1};
                  \mathsf{Pop}^{9-j}
\]
returns $(q_j,w,0,0)$. The selected $q_j$ is not one, so every subsequent
Pop uses the second case of (\ref{eq:ns-pop}). The last Pop consumes the
saved zero and clears the cache. For $u\le a$, the dominated head is
selected by
\[
 \mathsf{Discard};\mathsf{MP}_u;
 \mathsf{AX}_{b,a-u,c};S^9;\mathsf{Pop}.
\]
This produces $(\rho(a,b,c,u),w,0,0)$. From the zero top/cache store,
$\mathsf{MP}_u$ gives $(0,w,u,0)$ for every $u$: at $u=0$ the two branches
agree, and at $u>0$ the implication code is positive and differs from the
zero cache. A permitted gate uses
\[
 \mathsf{Discard};\mathsf{AX}_{A,D,H};\mathsf{Pop}^{10},
\]
returning $(g(H,D,A),w,0,0)$. Truth uses Discard alone.

Generalization at variable label $i$ is implemented by
\[
 \mathsf{Gen}_i=\mathsf{Push};\mathsf{MP}_i;S;
                    \mathsf{AX}_{i,0,0};S^{10}.
\]
From $(a,w,0,0)$, the first three pieces give $(0,w,i,a)$.
Equation (\ref{eq:ns-ten}) then gives $(U(i,a),w,0,0)$.
This word is total even when $a=1$: its first Shift tests zero, and all
later Shifts test only generated non-one heads. No hidden assumption that
an arbitrary formula is nonfalse occurs here.

\subsection{Modus ponens and the reserved zero}

The compilation invariant is: from every $(t,w,0,0)$, the word $C(p)$
for a numerical-calculus derivation $p$ either halts or returns
$(\operatorname{concl}(p),w,0,0)$.
Axioms and generalization have just been handled. Suppose the premises
are $a$ and $I(a,b)$. In the next display, $C(a)$ and $C(I(a,b))$
mean compilation of the given premise derivations, not a search for proofs
of those formulas. Use
\begin{equation}
 \mathsf{Discard};\mathsf{Push};C(I(a,b));\mathsf{Push};
 C(a);\mathsf{Push};S;S;\mathsf{MP}_b;\mathsf{Pop}.
 \label{eq:ns-mp-word}
\end{equation}
The initial Discard/Push reserves a zero below both recursive premise
computations. If neither recursive word halts, the store just before the
two Shifts is
\[
                  (0,[a,I(a,b),0;w],0,0).
\]
The first Shift returns $(0,[I(a,b),0;w],0,a)$. The second halts if
$a=1$; otherwise it gives $(a,[0;w],0,I(a,b))$. MP succeeds and returns
$(b,[0;w],b,0)$. The final Pop halts if $b=1$ and otherwise returns
$(b,w,0,0)$. The offset is cleared and the reserved zero is consumed
exactly once. This verifies every branch of the required halt-or-return
contract. At $a=b=w=0$ the proof uses literal Cantor projections and
$I(0,0)=0$; it does not appeal to distinct syntactic constructors.

Induction on the finite derivation establishes the compilation invariant.
The recursive compiler is primitive recursive: each derivation rule is
replaced by a fixed finite word together with recursive premise words and
numerical parameters already present in the coded derivation. PA proves
its correctness by induction on that coded derivation.

\subsection{A detector valid from every carried store}

Appending a plain Shift is insufficient: $(1,0,0,0)$ would continue to
zero. Instead define
\[
                 \mathsf{Detect}=A_0;S^{11}=\mathsf{Pop}.
\]
By (\ref{eq:ns-pop}), Detect halts from $(1,w,d,h)$ for every natural
$w,d,h$. Given a derivation $p$ of code one, the concrete command word
\[
                     A_0;C(p);\mathsf{Detect}
 \label{eq:ns-universal-word}
\]
halts from \emph{every} incoming store. The first AX clears offset and
cache while leaving some top and stack; the compiler works with arbitrary
such top and stack. It either halts already or returns current top one,
which Detect then detects. This argument does not assume that the
scheduler resets its store between unsuccessful command words.

The unchanged scheduler has fresh counter initially two and a total
Cantor command decoder whose positive continuation codes decrease.
Let $E(\epsilon)=0$ and
$E(c::w)=1+\operatorname{encodeCommand}(c,E(w))$ be the positive
continuation convention. The exact fresh-code forms are
\[
\begin{aligned}
 \operatorname{encodeCommand}(S,z)&=p(0,z),\\
 \operatorname{encodeCommand}(\mathsf{MP}_x,z)&=p(1,p(x,z)),\\
 \operatorname{encodeCommand}(\mathsf{AX}_{A,x,y},z)
                     &=p(A+2,p(y,p(x,z))).
\end{aligned}
\]
A fresh candidate uses its raw code directly. A positive continuation $n$
first decodes raw $n-1$, whose residual is at most $n-1<n$. Thus a fresh
candidate $n$ takes at most $n+1$ decoded commands before its continuation
is exhausted, unless it halts earlier.
If $z=E(A_0;C(p);\mathsf{Detect})$, the fresh candidates
\[
                  p(k+2,p(0,p(0,z)))\qquad(k\in\mathbb N)
\]
form a strictly increasing, unbounded family. Their decoded words prepend
one arbitrary AX to the all-store halting word. Thus omitting the finitely
many candidates below two causes no loss. PA verifies termination of each
command macro and decoding loop by the inherited natural-number ranking
arguments. For a given proof, an explicit sufficiently large prefix code
and the finite search up to it produce a finite halting source run.
Evaluation of each candidate is a bounded traversal of at most $n+1$
commands, each with a total primitive-recursive macro implementation;
halted configurations can be treated as absorbing. A bounded outer
iteration up to the chosen candidate code then constructs the complete
finite source trace. Code construction and finite trace concatenation are
primitive recursive, and PA proves their correctness by ordinary induction.

\subsection{Soundness and the calibration dependency}

For soundness, annotate the current top $t$, cached formula $h$, and every
live entry in the canonical numerical stack with a proof in the retained
calculus, or with the inherited PA proof of its arithmetic translation.
For $w>0$ its Cantor tail is strictly below $w$, so the canonical stack is
obtained by a primitive-recursive bounded recursion. A popped zero stack
supplies truth, and pushing zero onto zero introduces no live entry.

AX adds the same valid axiom heads, the generalization of the old cache,
and the saved old top. MP success takes the annotated formulas $t$ and
$I(t,b)=h$ to the annotated conclusion $b$, retains the stack and annotates
its zero cache by truth. MP failure outputs truth and retains the stack.
The numerical offset is not required to be a theorem. Shift transports
annotations; if it halts, its cache is code one and already has a proof.
These are uniform syntactic proof transformations, so PA proves that a
finite source halt yields a PA refutation. No reflection or uniform truth
predicate for PA is used.

The inherited logical-calibration dependency is the conclusion-preserving,
primitive-recursive translation between PA refutations and refutations in
the separate-head numerical calculus. Its arithmetic interpretation uses
finite binary membership and PA's finite-comprehension proof templates;
its completeness side uses the established reverse membership rewriting
and neighborhood-equality interfaces. The unchanged proof texts are
retained under \texttt{references/logical}. The present change neither
replaces those arguments by the repository's conditional iPRA endpoints
nor claims a new kernel-checked proof of them.

Combining that dependency with the new all-store compiler and annotation
argument gives, for the new source $R_{NS}$,
\[
 \mathrm{PA}\vdash
 \bigl(\operatorname{Halt}(R_{NS})\leftrightarrow
       \exists p\,\operatorname{Prf}_{\mathrm{PA}}
                           (p,\ulcorner0=1\urcorner)\bigr).
\]

\subsection{From the source to the exact tape table}

The complete 4,750 primitive outcomes are checked for arbitrary natural
register values, both decrement branches, an arbitrary gap parameter and
an arbitrary finite nonempty list of spare register values. An independent
literal-tape checker reconstructs the source by the explicit address
splice, compares two control-flow resolvers, rederives the four generic
list lemmas from the actual transition rows, and checks the exact complete
frames. Increments include their specified spare-list effects.

Literal blank execution reaches the complete source-25 frame in 57,768
steps, at native state 84 and head coordinate $-2$, with logical registers
$(2,0,0,0,0,0,0,0,0)$ and 26 spare registers of value one. Two different
finite tape representations compare the entire tape, including its zero
exterior; there is no earlier halt. The dense/native map conjugates every
row exactly.

Every source step has a terminating positive-length physical simulation,
with no premature physical halt. A finite halting source trace therefore
produces a physical halting trace by finite concatenation after startup.
Conversely, if the physical machine first halted after $T$ steps while
its first $T+1$ source steps were all nonhalting, their positive simulation
segments would concatenate into the same deterministic physical run for
more than $T$ nonhalting steps, a contradiction. This yields both
finite-trace implications in PA. Consequently, for the exact delivered
194-state table,
\[
 \boxed{\mathrm{PA}\vdash
 \bigl(\neg\operatorname{Halt}(\mathrm{PA194NS},\mathrm{blank})
                    \leftrightarrow\operatorname{Con}(\mathrm{PA})\bigr).}
\]
The new result is an ordinary mathematical PA-provability argument with
finite and symbolic executable certificates. These certificates are not
a new end-to-end Lean/Rocq theorem; actual standard-model nonhalting is
conditional on PA's consistency. No state minimality conclusion follows
from the unsuccessful smaller-machine searches.
